% Provenance (never printed): sources/book2/BRIEFS.md, g10-06. Equilibrium as a linear system,
% excess demand and supply, comparative statics, linked markets. Real numbers:
% figdata/grade-10/06-market-equilibrium/ (EIA world and US oil production; BLS pump price and
% CPI-U as the price, since the crude price has no approved source). Brenland and Kestria invented.
\chapter{Market Equilibrium}\label{ch:g10:market-equilibrium}

Each point of the opening chart is one year of the world oil market: how many barrels the world produced a day, and what
a gallon of gasoline cost an American driver, in today's money. The points are joined in the order of the years. Some
steps go up and to the right: more oil and a higher price, as from 2020 to 2021. Others go down and to the right: more oil
and a lower price, as from 2014 to 2015. Others go down and to the left, as in 2009 and in 2020, when the world produced
less oil and paid less for it.
% figdata: world_oil (2014 93.95, 4.46; 2015 96.65, 3.24; 2019 100.5, 3.23; 2020 94.01, 2.64; 2021 95.72, 3.53)

These directions are not random. Every point is where that year's demand met that year's supply, and the direction of each
step says which of the two curves moved. This chapter puts the two curves of the last two chapters together, solves for the
point where they meet, and works out how that point moves when either curve does.

\begin{openingchart}{The world oil market, 2005--2024: world production of oil and other liquid fuels, in millions of barrels a
day, against the price of regular gasoline in the United States, in 2024 dollars per gallon. Points are joined in the order of the
years. The pump price stands in for the price of crude oil, which it follows.}{US Energy Information Administration,
international energy statistics; US Bureau of Labor Statistics, average retail prices and CPI-U all items}
\begin{tikzpicture}
\begin{axis}[omrecord, width=0.86\linewidth, height=7.2cm, xmin=83, xmax=105, ymin=2.4, ymax=5.2,
  xlabel={world production (million barrels a day)}, ylabel={dollars per gallon (2024 prices)}]
\addplot[omDef, thick, mark=*, mark size=1.4pt] table[x=q, y=price] {figdata/grade-10/06-market-equilibrium/out/world_oil.csv};
\addplot[draw=none, nodes near coords, point meta=explicit symbolic,
  nodes near coords style={font=\scriptsize, anchor=south west, inner sep=1.5pt}]
  table[x=q, y=price, meta=year] {figdata/grade-10/06-market-equilibrium/out/world_labels.csv};
\end{axis}
\end{tikzpicture}
\end{openingchart}

\begin{recall}
From the grades 6--9 volume: at the \emph{equilibrium price} the quantity demanded equals the quantity supplied; below it there is a
\emph{shortage}, above it a \emph{surplus}, and the price tends to move towards equilibrium. From mathematics: two straight lines
with different slopes cross at exactly one point, found by solving a system of two linear equations. The grades 6--9
volume also solved $a - bp = c + dp$ for $p^* = (a - c)/(b + d)$, and showed that a shift of demand by $k$ moves the price by
$k/(b + d)$; this chapter restates both and builds on them. From \cref{ch:g10:demand} and
\cref{ch:g10:supply}: \omterm{def:g10:demand:market-demand}{market demand} slopes down, \omterm{def:g10:supply:market-supply}{market supply} slopes up, and each shifts when something other than the price changes.
\end{recall}

\section{Equilibrium as a pair of equations}\label{sec:g10:market-equilibrium:system}

The model of this chapter describes a market in which no buyer and no seller is large enough to move the price.

\begin{definition}[Price taker]\label{def:g10:market-equilibrium:price-taker}
A buyer or seller is a \emph{price taker}\index{price taker} when it buys or sells so small a share of the market that it takes the
market price as given: it can buy or sell as much as it wants at that price, and it cannot change the price by changing its own
quantity.
\end{definition}

Write the \omterm{def:g10:demand:market-demand}{market demand} and supply as straight lines, $Q^d = a - b\,p$ and $Q^s = c + d\,p$, with $b > 0$ and $d > 0$. The equilibrium
price $p^*$ is the price at which $Q^d = Q^s$.

\begin{proposition}[The equilibrium of a linear market, recalled]\label{prop:g10:market-equilibrium:solution}
If $b > 0$, $d > 0$ and $a > c$, the market has exactly one equilibrium,
\[
  p^* = \frac{a - c}{b + d}, \qquad Q^* = \frac{a\,d + b\,c}{b + d},
\]
and $p^* > 0$.
\end{proposition}
\begin{proof}
As in the grades 6--9 volume: $(b + d)\,p = a - c$ has exactly one solution since $b + d > 0$, positive because $a > c$; then
$Q^* = a - b\,p^* = (ad + bc)/(b + d)$.
\end{proof}

The condition $a > c$ says that at a price of zero, buyers want more than sellers offer; without it, the good would be free, like air.

\begin{example}[Bread in a Brenland town]\label{ex:g10:market-equilibrium:bread}
The \omterm{def:g10:demand:market-demand}{market demand} for bread is $Q^d = 50 - 3p$ and its supply $Q^s = 2p - 10$, in hundreds of loaves a day, with $p$ in dalens. Then $50 -
3p = 2p - 10$, so $5p = 60$ and $p^* = 12$ dalens; $Q^* = 50 - 36 = 14$ hundred loaves (check: $24 - 10 = 14$).
\end{example}

\section{Excess demand, excess supply and how the price moves}\label{sec:g10:market-equilibrium:excess}

\begin{definition}[Excess demand and excess supply]\label{def:g10:market-equilibrium:excess-demand}
At a price $p$, \emph{excess demand}\index{excess demand} is the quantity $Q^d(p) - Q^s(p)$ when it is positive; \emph{excess
supply}\index{excess supply} is the quantity $Q^s(p) - Q^d(p)$ when it is positive.
\end{definition}
\phantomsection\label{def:g10:market-equilibrium:excess-supply}

The shortage and surplus of the grades 6--9 volume are now measured: in the bread market at 10 dalens, $Q^d = 20$ and $Q^s = 10$, an
\omterm{def:g10:market-equilibrium:excess-demand}{excess demand} of 10 hundred loaves; at 14 dalens, $Q^d = 8$ and $Q^s = 18$, an \omterm{def:g10:market-equilibrium:excess-demand}{excess supply} of 10.

\begin{proposition}[Excess demand falls as the price rises]\label{prop:g10:market-equilibrium:excess}
In a linear market, $Q^d(p) - Q^s(p) = (a - c) - (b + d)\,p$. It is positive below $p^*$, zero at $p^*$ and negative above $p^*$, and
every dalen of price reduces it by $b + d$.
\end{proposition}
\begin{proof}
$Q^d - Q^s = a - bp - c - dp = (a - c) - (b + d)p$, a straight line of slope $-(b + d) < 0$ in $p$, equal to zero at $p = (a - c)/(b +
d) = p^*$. A decreasing line is positive before its zero and negative after.
\end{proof}

How does a market find $p^*$? Nobody computes it. When there is \omterm{def:g10:market-equilibrium:excess-demand}{excess demand}, buyers find the shelves empty and some offer more; sellers
notice that they sell out early and raise their price. When there is \omterm{def:g10:market-equilibrium:excess-demand}{excess supply}, unsold stock piles up and sellers cut prices to clear
it. Each step moves the price towards $p^*$, and \cref{prop:g10:market-equilibrium:excess} says why: on the side of $p^*$ where the price is,
the gap always pushes it the right way. How fast this happens differs from market to market: in an auction of fish it takes minutes, in the
market for flats it can take years.

\begin{omfigure}
\begin{tikzpicture}[econ, x=0.28cm, y=0.5cm, yshift=-2cm]
% vertical coordinate = p - 4 (the axis starts at a price of 4 dalens)
\draw[econ axes] (0,13.5) node[left, font=\small] {$p$} |- (31,0) node[below left=12pt and 0pt, font=\small] {$Q$ (hundreds of loaves)};
\draw[demand] (5,11) -- (29,3);
\draw[supply] (0,1) -- (24,13);
\draw[guide] (0,8) -- (14,8) -- (14,0);
\node[eqpoint] at (14,8) {};
\draw[guide] (0,6) -- (10,6);
\draw[|<->|, omThm!80!black] (10,6) -- (20,6);
\node[font=\footnotesize, anchor=north] at (15,5.8) {excess demand 10};
\draw[guide] (0,10) -- (8,10);
\draw[|<->|, omDef!80!black] (8,10) -- (18,10);
\node[font=\footnotesize, anchor=south] at (13,10.2) {excess supply 10};
\foreach \p in {6,8,10,12,14,16} \draw (0,\p-4) -- (-0.3,\p-4) node[left, font=\footnotesize] {\p};
\foreach \x in {8,10,14,18,20} \draw (\x,0) -- (\x,-0.2) node[below, font=\footnotesize] {\x};
\node[font=\footnotesize, text=omDef!80!black, anchor=north east] at (29,2.6) {$Q^d = 50 - 3p$};
\node[font=\footnotesize, text=omThm!80!black, anchor=north west] at (24.3,13) {$Q^s = 2p - 10$};
\end{tikzpicture}

{\small The bread market. At 12 dalens the market clears with 14 hundred loaves. At 10 dalens buyers want 20 and sellers offer 10; at 14
dalens sellers offer 18 and buyers want 8.}
\end{omfigure}

\section{Comparative statics}\label{sec:g10:market-equilibrium:statics}

\begin{definition}[Comparative statics]\label{def:g10:market-equilibrium:comparative-statics}
\emph{Comparative statics}\index{comparative statics} compares the equilibrium before and after a change in an \omterm{def:g10:economic-models:endogenous}{exogenous variable}, \omterm{def:g10:economic-models:other-things-equal}{other
things equal}, without describing the path from one to the other.
\end{definition}

\begin{proposition}[One shift; the price part recalled]\label{prop:g10:market-equilibrium:one-shift}
If demand shifts right by $\Delta > 0$ at every price ($a$ becomes $a + \Delta$), the equilibrium price rises by $\Delta/(b + d)$ and the
quantity by $d\,\Delta/(b + d)$. If supply shifts right by $\Sigma > 0$ ($c$ becomes $c + \Sigma$), the price falls by $\Sigma/(b + d)$ and
the quantity rises by $b\,\Sigma/(b + d)$.
\end{proposition}
\begin{proof}
By \cref{prop:g10:market-equilibrium:solution}, replacing $a$ by $a + \Delta$ changes $p^*$ by $\Delta/(b + d)$ and $Q^*$ by $d\Delta/(b+d)$.
Replacing $c$ by $c + \Sigma$ changes $p^*$ by $-\Sigma/(b + d)$ and $Q^*$ by $b\Sigma/(b + d)$.
\end{proof}

A rise in demand raises both the price and the quantity; a rise in supply raises the quantity and lowers the price. A fall in either does
the reverse. This gives a rule for reading data: when price and quantity move the same way, the demand shift dominated; when they move in
opposite ways, the supply shift dominated.

\begin{example}[More bread buyers]\label{ex:g10:market-equilibrium:more-buyers}
A new school opens in the town and the demand for bread becomes $Q^d = 60 - 3p$: $\Delta = 10$. The price rises by $10/5 = 2$ dalens, to 14,
and the quantity by $2 \times 10/5 = 4$, to 18 hundred loaves (check: $60 - 42 = 18$ and $28 - 10 = 18$). Of the 10 extra hundred loaves
wanted at the old price, only 4 are bought: the higher price discourages part of the new demand and draws out more supply.
\end{example}

When both curves move at once, the two effects add.

\begin{proposition}[Two shifts at once]\label{prop:g10:market-equilibrium:two-shifts}
If demand shifts right by $\Delta > 0$ and supply right by $\Sigma > 0$, the equilibrium quantity rises, and the price changes by
$(\Delta - \Sigma)/(b + d)$: it rises if $\Delta > \Sigma$, falls if $\Delta < \Sigma$, and does not change if $\Delta = \Sigma$.
\end{proposition}
\begin{proof}
Adding the two changes of \cref{prop:g10:market-equilibrium:one-shift}: the price changes by $\Delta/(b+d) - \Sigma/(b+d)$ and the quantity
by $(d\Delta + b\Sigma)/(b+d)$, which is positive since every term is.
\end{proof}

The sign of one of the two changes is always known and the other depends on the sizes of the shifts. The four cases are drawn below.

\begin{omfigure}
\begin{tikzpicture}[econ, x=0.42cm, y=0.38cm]
\foreach \col/\row/\dd/\ss/\title in {0/0/1/1/{demand $\uparrow$, supply $\uparrow$: $Q \uparrow$, $p$ ?},
                                       1/0/1/-1/{demand $\uparrow$, supply $\downarrow$: $p \uparrow$, $Q$ ?},
                                       0/1/-1/1/{demand $\downarrow$, supply $\uparrow$: $p \downarrow$, $Q$ ?},
                                       1/1/-1/-1/{demand $\downarrow$, supply $\downarrow$: $Q \downarrow$, $p$ ?}} {
  \begin{scope}[shift={(\col*16,-\row*12)}]
    \draw[econ axes] (0,9.5) node[left, font=\scriptsize] {$p$} |- (12,0) node[below, font=\scriptsize] {$Q$};
    \draw[demand] (1,8.5) -- (9,1.5);
    \draw[supply] (1,1.5) -- (9,8.5);
    \draw[demand, shifted] (1+2*\dd,8.5) -- (9+2*\dd,1.5);
    \draw[supply, shifted] (1+2*\ss,1.5) -- (9+2*\ss,8.5);
    \node[eqpoint] at (5,5) {};
    \node[font=\scriptsize, align=center] at (6,-1.6) {\title};
  \end{scope}}
\end{tikzpicture}

{\small The four cases of two shifts at once (dashed: the new curves). In each case one change is certain and the other depends on which
shift is larger.}
\end{omfigure}

The opening chart can now be read. From 2014 to 2015, world production rose from 93.95 to 96.65 million barrels a day while the price fell
from \usd{4.46} to \usd{3.24} a gallon: a supply-led move. The second chart shows where the supply came from: of the world's rise of 9.75
million barrels a day between 2008 and 2015, 6.59 million came from the United States, about two-thirds. From 2019 to 2020, production fell
from 100.50 to 94.01 million barrels a day and the price from \usd{3.23} to \usd{2.64}: a demand-led move, as travel stopped in the first
year of the pandemic. Of the 19 yearly steps of the chart, 13 moved price and quantity the same way. Demand moves this market most years;
the large supply moves are rarer.
% figdata: world_oil (2008 q 86.90, us 8.56; 2015 q 96.65, us 15.15); world_moves (13 of 19 same-direction)

\begin{omchart}{Production of oil and other liquid fuels, United States and the rest of the world, 2005--2024, in millions of barrels a
day. Both panels span 25 million barrels a day.}{US Energy Information Administration, international energy statistics}
\begin{tikzpicture}
\begin{groupplot}[group style={group size=2 by 1, horizontal sep=1.6cm},
  omtime, width=0.48\linewidth, height=5.4cm, xmin=2004, xmax=2025, xtick={2005,2010,...,2020}]
\nextgroupplot[width=0.48\linewidth, title={\small United States}, ymin=0, ymax=25]
\addplot[omThm, thick] table[x=year, y=us] {figdata/grade-10/06-market-equilibrium/out/world_oil.csv};
\nextgroupplot[width=0.48\linewidth, title={\small rest of the world}, ymin=60, ymax=85]
\addplot[omDef, thick] table[x=year, y=row] {figdata/grade-10/06-market-equilibrium/out/world_oil.csv};
\end{groupplot}
\end{tikzpicture}
\end{omchart}

\section{Linked markets}\label{sec:g10:market-equilibrium:linked}

Markets do not stand alone. When the price of one good changes, the demand for its \omterm{def:g10:demand:substitutes}{substitutes} and \omterm{def:g10:demand:substitutes}{complements} shifts, and the supply of
goods that use it as an input shifts. A shock in one market spreads.

\begin{example}[A frost and two fruits]\label{ex:g10:market-equilibrium:frost}
A spring frost destroys part of the apple blossom. In the apple market, supply shifts left: the price of apples rises and fewer are sold.
Pears are a substitute for apples, so the dearer apples shift the demand for pears right: the price of pears rises and more pears are sold,
although nothing happened to pear trees. Apple juice uses apples as an input: its supply shifts left. And the dearer pears in turn shift the
demand for apples a little to the right, which raises the price of apples a little further. The weekend problem computes all of this.
\end{example}

\begin{omfigure}
\begin{tikzpicture}[econ, x=0.42cm, y=0.38cm]
\begin{scope}
\draw[econ axes] (0,9.5) node[left, font=\small] {$p$} |- (12,0) node[below, font=\small] {apples};
\draw[demand] (1,8.5) -- (9,1.5);
\draw[supply] (1,1.5) -- (9,8.5);
\draw[supply, shifted] (0.3,3.08) -- (6.5,8.5);
\node[eqpoint] at (5,5) {};
\node[eqpoint, fill=omThm] at (3.75,6.09) {};
\draw[phase arrow] (7.1,7) -- (5.0,7);
\node[font=\scriptsize, anchor=west] at (7.3,7) {frost};
\end{scope}
\begin{scope}[xshift=7.4cm]
\draw[econ axes] (0,9.5) node[left, font=\small] {$p$} |- (12,0) node[below, font=\small] {pears};
\draw[demand] (1,8.5) -- (9,1.5);
\draw[supply] (1,1.5) -- (9,8.5);
\draw[demand, shifted] (2.5,8.5) -- (10.5,1.5);
\node[eqpoint] at (5,5) {};
\node[eqpoint, fill=omDef] at (5.75,5.53) {};
\draw[phase arrow] (5.7,4.5) -- (6.95,4.5);
\node[font=\scriptsize, anchor=west] at (7.4,4.5) {dearer apples};
\end{scope}
\end{tikzpicture}

{\small A frost shifts the supply of apples left (left panel): the apple price rises and fewer apples are sold. Dearer apples shift the
demand for pears, a substitute, right (right panel): the pear price and quantity rise.}
\end{omfigure}

\begin{model}{The competitive market}
\assumes Many price-taking buyers and sellers of one good of known quality; the price is free to move and moves towards the level at which
  the quantities demanded and supplied are equal.
\predicts The equilibrium price and quantity; after one shift, the direction of both and their size, from the slopes; after two shifts, the
  direction of one of them.
\tested Markets run in laboratories, with buyers and sellers given private values and costs, trade after a few rounds at prices and
  quantities close to the predicted equilibrium (Smith 1962). Consistent with it: the directions of the world oil market's yearly steps fit
  the shifts known from outside the market, a supply-led step in 2015 after the rise of US output, demand-led falls in the recession of 2009
  and in 2020, a demand-led rise in 2021.% ledger: g10:smith-1962

\breaks When there are few buyers or sellers (grade 11), when a rule fixes the price (\cref{ch:g10:what-prices-do}), when prices move
  slowly (the market for flats, wages), and when buyers cannot see the quality of the good (grade 11).
\end{model}

\begin{record}[The oil market, 2014--2016 and 2020--2022]
% figdata: world_oil (2014 93.95, 4.46; 2016 97.09, 2.80; 2020 94.01, 2.64; 2022 100.47, 4.39)
Between 2014 and 2016, world production rose from 93.95 to 97.09 million barrels a day and the price at the pump fell from \usd{4.46} to
\usd{2.80}. Supply had shifted right: American output, which had grown by 5.6 million barrels a day since 2008 with the new wells of
\cref{ch:g10:supply}, rose by 0.7 million more, and other producers raised theirs by 2.5 million as the price fell. The model predicts what
was seen, more oil at a lower price.
% figdata: world_oil (us 2008 8.56, 2014 14.16, 2016 14.85; row 2014 79.79, 2016 82.24) Between 2020 and 2022, production rose from 94.01 to 100.47 million barrels a day and the
price from \usd{2.64} to \usd{4.39}: demand recovered after the first year of the pandemic, and the model predicts a rise of both. The record
also shows the model's limit: it says in which direction each curve moved, given the directions of price and quantity, but it cannot say by
how much each moved when both moved at once, as they surely did in 2022.
\end{record}

\section{Exercises}

\begin{exercise}[$\star$]\label{exo:g10:market-equilibrium:1}
Find the equilibrium price and quantity when $Q^d = 120 - 4p$ and $Q^s = 6p - 30$.
\end{exercise}

\begin{exercise}[$\star$]\label{exo:g10:market-equilibrium:2}
Find the equilibrium when $Q^d = 90 - p$ and $Q^s = 2p$. Check your answer in both equations.
\end{exercise}

\begin{exercise}[$\star$]\label{exo:g10:market-equilibrium:3}
In the bread market of \cref{ex:g10:market-equilibrium:bread}, compute the \omterm{def:g10:market-equilibrium:excess-demand}{excess demand} or \omterm{def:g10:market-equilibrium:excess-demand}{excess supply} at 8, 11 and 13 dalens.
\end{exercise}

\begin{exercise}[$\star$]\label{exo:g10:market-equilibrium:4}
Read the opening chart. In which two years was the price highest? Between 2016 and 2018, did the price and the quantity move the same way?
\end{exercise}

\begin{exercise}[$\star$]\label{exo:g10:market-equilibrium:5}
For each change, say which curve of the market for hotel rooms in a seaside town shifts, and in which direction: (a) a new hotel opens;
(b) a heat wave inland; (c) the wages of hotel staff rise; (d) a new motorway halves the journey from the capital.
\end{exercise}

\begin{exercise}[$\star\star$]\label{exo:g10:market-equilibrium:6}
In the bread market, a rise in the price of flour shifts supply to $Q^s = 2p - 20$. Find the new equilibrium, and check it against
\cref{prop:g10:market-equilibrium:one-shift}.
\end{exercise}

\begin{exercise}[$\star\star$]\label{exo:g10:market-equilibrium:7}
In the bread market, demand shifts right by 15 and supply right by 5 at the same time. Use \cref{prop:g10:market-equilibrium:two-shifts} to
find the changes in price and quantity, then check by solving the new system.
\end{exercise}

\begin{exercise}[$\star\star$]\label{exo:g10:market-equilibrium:8}
Coffee and tea are \omterm{def:g10:demand:substitutes}{substitutes}. A disease destroys part of the coffee harvest. Predict the direction of the change in the price and quantity of
coffee, and in the price and quantity of tea. Draw the two markets.
\end{exercise}

\begin{exercise}[$\star\star$]\label{exo:g10:market-equilibrium:9}
Read the second chart. By how many million barrels a day did US production rise between 2008 and 2015, and by how many the rest of the world's?
What share of the world's rise came from the United States?
\end{exercise}

\begin{exercise}[$\star\star$]\label{exo:g10:market-equilibrium:10}
The demand for a good is $Q^d = 48 - 2p$. Three sellers have the supplies $q_1 = p$, $q_2 = p - 4$ (for $p \ge 4$) and $q_3 = 2p - 20$ (for $p \ge
10$). Find the equilibrium price and quantity. (Try first the piece of the \omterm{def:g10:supply:market-supply}{market supply} on which it lies.)
\end{exercise}

\begin{exercise}[$\star\star$]\label{exo:g10:market-equilibrium:11}
In \cref{ex:g10:market-equilibrium:more-buyers}, 10 hundred loaves more were wanted at the old price, but only 4 more were bought. Explain how the
\omterm{def:g10:market-equilibrium:excess-demand}{excess demand} of 10 at the old price is closed: how much by buyers who give up because of the higher price, and how much by sellers who offer
more?
\end{exercise}

\begin{exercise}[$\star\star\star$]\label{exo:g10:market-equilibrium:12}
Over one year, the price of a good rose and the quantity sold rose. Someone concludes that ``supply did not change''. Explain why this is not a
valid conclusion, and what the observation does show.
\end{exercise}

\begin{exercise}[$\star\star\star$]\label{exo:g10:market-equilibrium:13}
Prove that if demand shifts left by $\Delta$ and supply shifts right by $\Sigma$, the price falls, and find the condition under which the quantity
rises.
\end{exercise}

\begin{exercise}[$\star\star\star$]\label{exo:g10:market-equilibrium:14}
In the bread market, the town fixes the price of bread at 10 dalens and forbids any other price. How many loaves are sold? Who gets them, if bakers
sell to the first customers to arrive? What would you expect to see in the shops in the morning?
\end{exercise}

\begin{exercise}[$\star\star\star$]\label{exo:g10:market-equilibrium:15}
Suppose that, in some market, the supply curve sloped down, $Q^s = c - d\,p$ with $d > 0$, as steeply as or more steeply than demand. Show that
if $d = b$ the system $Q^d = Q^s$ has no solution or infinitely many, and that if $d > b$ and $a > c$ its only solution is a negative
price. Which step of the proof of \cref{prop:g10:market-equilibrium:solution} fails?
\end{exercise}

\section{Problem: A Frost in Kestria's Orchards}

\begin{problem}[{Weekend problem --- apples, pears and a four-unknown system}]
\label{pb:g10:market-equilibrium:1}
In Kestria, apples and pears are \omterm{def:g10:demand:substitutes}{substitutes}. With $p_A$ and $p_P$ the prices of apples and pears in kests per kilogram, and quantities in
tonnes a week, the markets are:
\[
\begin{aligned}
  \text{apples:}\quad & Q^d_A = 100 - 2p_A + p_P, & Q^s_A &= 3p_A - 20, \\
  \text{pears:}\quad  & Q^d_P = 80 - 2p_P + p_A,  & Q^s_P &= 2p_P - 10 .
\end{aligned}
\]

\medskip
\noindent\textbf{Part I --- The two markets before the frost.}
\begin{enumerate}
\item What happens to the demand for pears when the price of apples rises by one kest? Why does this make them \omterm{def:g10:demand:substitutes}{substitutes}?
\item Write the apple market's equilibrium condition as an equation in $p_A$ and $p_P$.
\item Do the same for the pear market.
\item Solve the system: find $p_A$ and $p_P$.
\item Find the quantities of apples and pears sold.
\end{enumerate}

\medskip
\noindent\textbf{Part II --- The frost.}
\begin{enumerate}[resume]
\item A frost destroys blossom, and the supply of apples becomes $Q^s_A = 3p_A - 58$. By how much has the supply of apples fallen at every price?
\item If the price of pears stayed at its old value, what would the new apple price and quantity be?
\item At that apple price, what would happen in the pear market at the old pear price: \omterm{def:g10:market-equilibrium:excess-demand}{excess demand} or \omterm{def:g10:market-equilibrium:excess-demand}{excess supply}, and how much?
\end{enumerate}

\medskip
\noindent\textbf{Part III --- The new equilibrium.}
\begin{enumerate}[resume]
\item Write the two equilibrium conditions after the frost.
\item Solve for the new prices.
\item Find the new quantities, and check each in its demand and its supply.
\item By how much did the price of apples rise? How much of the rise is due to the pear market's response?
\item The frost did not touch a single pear tree. Explain why pear growers sold more, at a higher price.
\end{enumerate}

\medskip
\noindent\textbf{Part IV --- What the shoppers and growers see.}
\begin{enumerate}[resume]
\item Compute the apple growers' takings before and after the frost (price times quantity, in kests per week, with prices per kilogram and
  quantities in tonnes).
\item Compute the pear growers' takings before and after.
\item Which group of growers gained from the frost? Can apple growers as a whole gain from a frost? What does it depend on?
\item By what percentage did each price rise, and each quantity change?
\item In the apple market, was the fall in the quantity sold a move along the demand curve, a shift of it, or both? And in the pear market?
\item Suppose instead that apples and pears were not linked (the price of each fruit did not enter the demand for the other) and that, before the
  frost, the two markets had the same equilibria as in question 4. What would the frost have done to the price of pears?
\item State the result of the problem: the new equilibrium prices and quantities in both markets, compared with those before the frost.
\end{enumerate}
\end{problem}
